\section{Computation and Universality}

Before any talk of life or mind, one fact licenses the whole climb. The substrate can compute anything. This is the permission slip for the entire theory. The wild claim is that one law on a crystal yields atoms, weather, and minds. It rests on a modest fact. The knit, the collide-then-stream law that drives the mesh, is \emph{universal}. Any pattern that can exist at all is a pattern the knit can host.

\begin{principle}[A simple law can carry any world]
A complex world does not require a complex law. The richness of what unfolds is not bought by complicating the rule. It is bought by the rule being universal, so that every expressible pattern is already within reach.
\end{principle}

Two terms used throughout deserve a plain statement first. A \emph{universal} computer is one that can simulate any other computer, so that whatever is computable at all is computable by it. The standard reference machine is the Turing machine, and a system is called Turing-universal when it can simulate one. The \emph{knit} is the state-law of the mesh, collide then stream. Each beat the tones in a dock interact by a fixed reversible charge-conserving table, then each tone moves one dock along its site. A \emph{dock} is one cell of the mesh, a here-now holding $24$ sites, and a \emph{site} is one of the $24$ directions out of a dock, each carrying a single tone in $\{-1, 0, +1\}$.

\subsection{The substrate is a computer}

The flow is a cellular automaton on the nose, not by analogy. It is a discrete-state, $24$-neighbour automaton on the $\{3,4,3,4\}$ mesh. Each dock holds $24$ tones, one per site, each tone in $\{-1, 0, +1\}$. There is no outside machine running it. No operator, no screen, no interpreter standing apart from the world. Running the knit on the mesh is what existence does. The computation and the existing are the same act.

Universality stands on four legs, all verified in code. The first three give \emph{weak universality}, the standard notion for a hyperbolic grid, where the computation lives against an infinite or ultimately periodic background. The fourth gives \emph{strong universality}, computation on the flat finite space the theory's physics actually inhabits.

\begin{table}[ht]
\centering
\begin{tabular}{@{}clll@{}}
\toprule
leg & claim & what it gives & status \\
\midrule
1 & the addressing tree is present & the Margenstern structure on $\{3,4,3,4\}$ & verified \\
2 & the knit is functionally complete & NAND, then Rule 110 & verified \\
3 & a conserving register machine & a two-counter Minsky machine & verified \\
4 & strong universality from the cusp & Conway's Life on the flat 3D space & verified \\
\bottomrule
\end{tabular}
\caption{The four legs of universality, weak from the first three, strong from the fourth.}
\label{tab:universality-legs}
\end{table}

\subsubsection{Leg one, the addressing tree}

Margenstern proved the hyperbolic grids $\{7,3\}$, $\{5,4\}$, and $\{5,3,4\}$ universal using a fixed kit. A Fibonacci addressing tree, a small set of son types, a preferred son, and a railway model of a train running on switched tracks. Every piece of that kit has a $\{3,4,3,4\}$ analog.

\begin{itemize}
\item the \emph{addressing tree}, with growth ratio about $18.4$, a bushier generalized-Fibonacci tree
\item the \emph{splitting matrix} whose region-types play the role of the son-types
\item the \emph{digit-zero preferred son} as the spine that runs straight outward
\item the full \emph{$24$ interior directions}, abundant room for railway switches and crossings, since a junction needs only $3$ edge-disjoint directions
\end{itemize}

The kit transfers wholesale. Whatever Margenstern's railway needs, the $24$-site dock supplies and more.

\subsubsection{Leg two, the knit is functionally complete}

The model's own signed-majority update, with a $+1$ bias and two $-1$ fills, computes NAND. NAND is functionally complete, the single gate from which every Boolean function is built. So the knit builds NOT, AND, OR, XOR, and a full adder, each as a fixed local arrangement of tones. In particular it builds Rule 110, the elementary cellular automaton that Cook proved universal. The knit-built Rule 110 reproduces a reference Rule 110 exactly, cell for cell, across the run.

\begin{proposition}[The knit is Cook-universal]
The signed-majority collision of the knit computes NAND. NAND is functionally complete, so the knit realizes every Boolean gate and in particular Rule 110, which is universal. The substrate's own ternary knit therefore hosts a Cook-universal automaton.
\end{proposition}

The tones are $-1$, $0$, $+1$, read as pain, peace, pleasure. The same three values the knit reads to update a dock are the values that carry the logic. There is no separate logical alphabet bolted on. The vibe itself is the bit, the trit rather, and the law that moves the world is the law that runs the gate.

\subsubsection{Leg three, a conserving register machine}

A two-counter Minsky machine is Turing-universal, and it runs directly on the mesh. The register machine is not a circuit laid over the substrate. It is the substrate's own dynamics arranged to count.

\begin{table}[ht]
\centering
\begin{tabular}{@{}ll@{}}
\toprule
machine part & substrate realization \\
\midrule
register & a charged subtree, navigated by address \\
increment & the lean creates a balanced pair \\
decrement & annihilation \\
test-zero & perception counts the charges \\
\bottomrule
\end{tabular}
\caption{A two-counter Minsky machine realized in the substrate's own conserving dynamics.}
\label{tab:minsky}
\end{table}

Total tone stays exactly conserved throughout. This is genuine substrate dynamics, not a logic gate dressed up. It runs real programs. The verified runs include ADD ($3 + 4 = 7$) and MULTIPLY ($3 \times 4 = 12$, $5 \times 5 = 25$), with the net tone conserved at every beat. The increment is the lean leaning out a balanced pair, the decrement is that pair's annihilation, and the test-for-zero is the same perception the knit already performs.

\subsubsection{Leg four, strong universality from the cusp}

The cusp of the $\{3,4,3,4\}$ mesh is the Euclidean cubic honeycomb $\{4,3,4\}$, which is $\mathbb{Z}^3$, the flat three-dimensional space the theory's physics lives in. Conway's Game of Life on a $\mathbb{Z}^2$ layer of it is a known universal automaton.

This is run, not merely argued. A glider seeded on the actual cusp plane evolves identically to a reference $\mathbb{Z}^2$ Life and translates by $(1, -1)$ after its period-$4$ cycle, exactly as a Conway glider must.

\begin{result}[The emergent space is itself a universal computer]
A Conway glider seeded on the cusp plane of $\{3,4,3,4\}$ reproduces reference $\mathbb{Z}^2$ Life step for step. The theory's emergent physical space, the flat $3$D space where gravity and the light cone live, can compute anything computable.
\end{result}

\subsection{Capacity, not outcome}

Universality is a statement about what can be expressed, not about freedom of result. The capacity is unlimited. The outcome of any given run is one. The fire can take any shape, but a particular fire takes one shape, and the only way to see which shape it takes is to let it burn. Universality says the world is not limited in advance to a small menu of behaviours. It does not say the world is undetermined, only that its determination is not narrow.

\subsection{Undecidability comes with the power}

The moment a system can simulate a universal machine, its halting becomes undecidable. This is not a flaw. It is the formal signature of a world whose future can only be known by living it. The knit still produces a perfectly definite future, beat by beat, with no ambiguity. There is simply no algorithm that fast-forwards to the answer in general. This is the same computational irreducibility that underwrites the freedom argument elsewhere in the paper.

\begin{theorem}[Undecidability is native to the geometry]
A universal substrate has an undecidable halting problem. On $\{3,4,3,4\}$ the natively undecidable problem is the domino problem, whether a given set of tile types can cover the plane. It is undecidable here because the substrate is a tiling. Undecidability is therefore a feature of the geometry itself, not an import from outside.
\end{theorem}

The future is fixed and yet unforeseeable by shortcut. That is the precise shape of a lived world. The determinism is total and the irreducibility is total, and the two hold at once.

\subsection{Super-Turing, a clearly speculative aside}

If a dock had infinitely many neighbours, the so-called infinigrid, it could evaluate a quantifier in a single beat and decide every level of the arithmetical hierarchy. This is hypercomputation, computation beyond the Turing limit. We keep it clearly labelled as speculation, because it rests on objects the theory does not have.

\begin{itemize}
\item it needs unbuildable infinite objects, a dock of infinite degree
\item the $\{3,4,3,4\}$ mesh has $24$-direction docks, not infinitely many, so the result does not apply to the actual substrate
\end{itemize}

The honest division is sharp. A clean theorem on undecidability for the real substrate, and a clearly flagged thought experiment on hypercomputation for a substrate we do not have. The two are never confused.

\subsection{The addressing idea in its cleanest form}

For the universe to be computable and navigable, each dock must know where it is and who its neighbours are from local information alone. The whole scheme reduces to a single idea. A dock's address is a word, and the word is the route to it.

The idea is most transparent on the two-dimensional face, the $\{7,3\}$ tiling. Number the docks shell by shell out from the center. Write each number in a \emph{Fibonacci numeration}, whose place values are $1, 2, 3, 5, 8, 13$, and so on. Every whole number is a unique sum of non-adjacent Fibonacci numbers, so no valid address ever has two ones in a row. The valid addresses form a clean language that a tiny automaton can recognize. The growth of the tree and the alphabet of the addresses are the same fact, seen twice.

The tree branches by \emph{sons}. Each dock has a small \emph{type}, which is just how many children it spawns, and the types reproduce one another by a fixed rule. On the two-dimensional face there are two types, a lean son and a full son. One \emph{preferred son} carries the spine straight outward, and the rest fan to the sides. That single recurrence grows the entire bulk.

Routing is then word surgery, nothing more.

\begin{itemize}
\item the \emph{parent} of a dock is its address with the last digit stripped off
\item a \emph{child} is its address with one more legal digit appended
\item to route from one dock to another, line up the two addresses from the left, walk up to the shared prefix by dropping digits, then walk down to the target by appending its remaining digits
\end{itemize}

The shared prefix names the common ancestor. The two tails name the two halves of the path. The whole route is read straight off the two names, with no map and no coordinates consulted.

\begin{result}[Routing is address arithmetic]
The route between any two docks is the difference of their two names. It is pure address arithmetic, computed in time about the logarithm of the distance, with no global map.
\end{result}

The $\{3,4,3,4\}$ mesh carries the very same idea, only richer. The growth root is near $18.3$ rather than the golden ratio, so the digit alphabet is the $24$-way child index in place of the Fibonacci binary. The route-from-the-root, the parent-by-strip, and the common-prefix walk are unchanged. There is one clean simplification on the way up. The two-dimensional face has \emph{cousins}, docks that sit on the same shell across different branches, and reconstructing them needs a second harder automaton. The $\{3,4,3,4\}$ mesh has no same-shell cousins. Every one of a dock's $24$ neighbours is a parent or a child. So that hard piece is simply absent here, and the richer tree is the simpler one to navigate.

\subsection{Addressing the mesh at scale}

For the universe to be computable and navigable in practice, every dock must compute its own neighbours from its own name, with no master map. Margenstern's splitting and Fibonacci-addressing program does this for $\{7,3\}$ and $\{5,3,4\}$. It has now been extended to $\{3,4,3,4\}$ and verified, measured to $2{,}000{,}000$ docks.

\begin{table}[ht]
\centering
\begin{tabular}{@{}ll@{}}
\toprule
result & measurement \\
\midrule
unique self-describing $O(\log n)$ address & $0$ duplicate addresses, $0$ decode round-trip failures \\
neighbour set from name alone & $1632/1633$ complete docks, the lone miss being the root \\
greedy address-routing & $1999/2000$ delivered, mean stretch $1.004$ \\
confluence automaton deterministic at window $K=2$ & $445/445$ confluence partners reconstructed \\
core adjacency invariants survive at scale & confirmed to $2{,}000{,}000$ docks \\
\bottomrule
\end{tabular}
\caption{Addressing $\{3,4,3,4\}$, measured to two million docks.}
\label{tab:addressing-scale}
\end{table}

The surprise that makes it clean is that $\{3,4,3,4\}$ has no same-shell cousins, so the hard two-dimensional piece of Margenstern's method is simply absent. The neighbour map is a finite generator-address transducer, a small automaton.

\begin{result}[Local self-knowledge]
A dock knows where it is and who its $24$ neighbours are from its name alone, by a finite automaton with no global map. The limit at two million docks is the coordinate-precision and memory wall, not the addressing, which is combinatorial and unbounded.
\end{result}

\subsubsection{Margenstern's sons, ported}

Margenstern's two-dimensional method gives each tree node a \emph{type}. The pentagrid has two, a black son with $2$ children and a white son with $3$, and the types reproduce each other by a fixed rule, with a preferred son continuing the spine. The $\{3,4,3,4\}$ port keeps the whole structure, richer in count but identical in shape.

\begin{table}[ht]
\centering
\begin{tabular}{@{}ll@{}}
\toprule
2D pentagrid & $\{3,4,3,4\}$ port \\
\midrule
$2$ son types (black, white) & $10$ son types (child counts $15$ to $24$) \\
$2 \times 2$ son-matrix, golden-ratio growth & $10 \times 10$ splitting matrix, Perron eigenvalue $18.284$ \\
preferred son & the digit-$0$ child, the spine successor \\
cousin automaton (hard) & none, no same-shell cousins \\
no non-tree structure & a small confluence transducer, deterministic at window $2$ \\
\bottomrule
\end{tabular}
\caption{Margenstern's son-types ported from the pentagrid to $\{3,4,3,4\}$.}
\label{tab:sons-ported}
\end{table}

The splitting matrix dominant eigenvalue, $18.284$, equals the measured shell growth ratio. The shells run $1, 24, 456, 8376$, with the ratio approaching $18.37$, and the characteristic polynomial of the matrix is the shell recurrence itself. So the $\{3,4,3,4\}$ tree is structurally simpler than the two-dimensional originals it descends from, despite being bushier.

\subsubsection{A worked address example}

Addresses are dot-separated digits. Each digit is the index of a dock among its parent's children, ordered deterministically by embedding, and the address is the route from the root. The root is the empty address, written \texttt{e}.

\begin{table}[ht]
\centering
\begin{tabular}{@{}rrlrr@{}}
\toprule
dock & shell & address & parent & children \\
\midrule
0 & 0 & \texttt{e} & -- & 24 \\
1 & 1 & \texttt{0} & 0 & 23 \\
24 & 1 & \texttt{23} & 0 & 15 \\
25 & 2 & \texttt{0.0} & 1 & 22 \\
26 & 2 & \texttt{0.1} & 1 & 21 \\
\bottomrule
\end{tabular}
\caption{A small slice of the addressing tree, with each dock's shell, address, parent, and child count.}
\label{tab:address-example}
\end{table}

A routing walk from dock $A$ to dock $B$ is pure address arithmetic. Compare the two addresses from the left, walk up from $A$ to the common ancestor by dropping the last digit, then walk down to $B$ by appending $B$'s remaining digits. Take $A$ as dock $481$, address \texttt{0.0.0}, and $B$ as dock $503$, address \texttt{0.1.0}. Their longest shared prefix is \texttt{0}, so the common ancestor is dock $1$.

\begin{table}[ht]
\centering
\begin{tabular}{@{}rrll@{}}
\toprule
step & dock & address & operation \\
\midrule
0 & 481 & \texttt{0.0.0} & start ($A$) \\
1 & 25 & \texttt{0.0} & drop last digit, go to parent \\
2 & 1 & \texttt{0} & drop last digit, common ancestor \\
3 & 26 & \texttt{0.1} & append digit 1, go to child \\
4 & 503 & \texttt{0.1.0} & append digit 0, reached $B$ \\
\bottomrule
\end{tabular}
\caption{A four-hop route read straight off the two addresses.}
\label{tab:route-walk}
\end{table}

Four hops, which equals the shortest-path distance. No map, no coordinates, only the two names.

\subsection{One structure, computation and physics}

Here is the decisive point. The computation and the physics sit on the same structure. Margenstern's heptagrid and pentagrid carry the Fibonacci addressing and the universality but nothing else. They are two-dimensional. No spinor, no Standard Model, no three-dimensional space. They are computation substrates and nothing more. The $\{3,4,3,4\}$ mesh carries the whole package and the physics.

\begin{itemize}
\item the same $24$ sites that give the bushy generalized-Fibonacci tree are the $D_4$ root system that gives the spinor and the three generations
\item the same cusp that runs Conway's Life is the flat three-dimensional space where gravity and the light cone live
\end{itemize}

\begin{result}[Navigation, universality, and physics on one mesh]
$\{3,4,3,4\}$ is the unique grid we have found where navigation, universality, and physics all sit on one structure. The old trade-off, that $\{7,3\}$ wins universality while $\{3,4,3,4\}$ wins spin and the two are complementary, is dissolved. $\{3,4,3,4\}$ has both, and is simpler, with no cousins, than the two-dimensional grids it descends from.
\end{result}

So the substrate is at once Lorentz-safe, exactly addressable, computationally universal (weakly through the ternary NAND, Rule 110, and the register machine, strongly through cusp Life), and physically loaded. One honest caveat belongs here. Universality covers the structure of any process, not the felt interior. The qualia, the light of experience, is taken as the base. That is a stance, not a derivation.

\subsection{Why it matters}

Computation is the floor under the whole ladder. Three facts hold together.

\begin{itemize}
\item the substrate can host any pattern, by universality
\item it cannot in general be shortcut, by undecidability
\item it is navigable from local names alone, by addressing
\end{itemize}

In the carrying image, the fire can take any shape, but the only way to see which shape it takes is to let it burn. This section establishes the capacity. The structures that capacity stores and recalls are taken up next, in the bulk as a data structure, the associative memory engine that looks up by content, and the data structures that recur across every tessellation.
